% !PRINTED 135
\weberpage{135}
%<<CONT

hence by (8) and (12)
\begin{equation}
 |\psi(x)|\,|x-\alpha|^a<|\psi_1(x)|\,|x-\alpha_1|^{a_1}\,|x-\alpha_2|^{a_2}\cdots+|\varphi(x)|.
\tag{16}
\end{equation}
If we now assume $x$ on the periphery of $\varrho$, and therefore put
\[
 |x-\alpha|=\varrho,
\]
then
\[
 |x-\alpha_1|<2\varrho,\qquad |x-\alpha_2|<2\varrho,\ldots,
\]
and consequently, by (10) and (16),
\begin{equation}
 \varrho^a A < (2\varrho)^{a_1+a_2+\cdots}B+|\varphi(x)|.
\tag{17}
\end{equation}
We now choose the upper bounds for the coefficients $\varepsilon_1,\varepsilon_2,
\ldots$ of $\varphi$ so small that
\[
 |\varphi(x)|<\varrho^aA'
\]
where $A'$ denotes a number which is smaller than $A$; then from (17) it follows that\begingroup\renewcommand{\thefootnote}{E1}\footnote[1]{[Source note:] Set $s=a_1+a_2+\cdots$. From (17) one obtains $A-A'<2^s\varrho^{s-a}B$. The contradiction concerns $s>a$, not $a>s$: for $0<\varrho<1$, one has $2^s\varrho^{s-a}B\le 2^nB\varrho\to0$. Thus the signs in (19), (20) should be $\ge$; summing still gives equality. The source reference (20), (21) means (19), (20). Printed text and formulas are retained.}\endgroup
\begin{equation}
 A-A'<2^{a}\varrho^{-a+a_1+a_2+\cdots}B.
\tag{18}
\end{equation}
But, for sufficiently small $\varrho$, this would no longer be possible if $a$ were greater than the sum $a_1+a_2+\cdots$. Hence
\begin{equation}
 a\le a_1+a_2+\cdots .
\tag{19}
\end{equation}
In the same way one proves that
\begin{equation}
 b\le b_1+b_2+\cdots,
 \qquad
 c\le c_1+c_2+\cdots,
 \qquad \ldots .
\tag{20}
\end{equation}
But since the sums of the left sides as well as of the right sides in the inequalities (20), (21) must be equal to $n$,
% [sic] Source-original mathematical cross-reference: the printed page says (20), (21), although the displayed inequalities are numbered (19), (20).
only equality signs can hold, and therefore
\[
 a=a_1+a_2+\cdots,
 \quad b=b_1+b_2+\cdots,
 \quad c=c_1+c_2+\cdots,
 \quad \ldots,
\]
whereby the second part of our theorem is also proved.

We shall add to the theorem just proved the following remark concerning the case of multiple roots.

The necessary and sufficient condition that $f(x)=0$ should have multiple roots at all is that $f(x)$ and $f'(x)$ should have a common factor. Therefore, if the algorithm
%CONT>>
